\section{讨论与结论}

三年数据实验表明，在区域风电集群日内预测中，未来气象仍是主导驱动，而预测起点前真实功率状态能够提供有效的运行状态约束。本文以该状态为条件读取气象记忆，提升了日内平均预测精度，并显著降低预测初期误差。实验结果进一步验证了显式信息角色划分和记忆层条件化在协调预测起点状态与未来气象信息方面的作用。纯气象模型、输出层融合、随机/零功率控制和气象侧消融为该机制及多尺度状态编码的附加作用提供了互补证据。当前证据仍限于单一区域的确定性预测，未来将进一步研究气象预报不确定性、跨区域迁移、模型简化和极端爬坡事件。
